\documentclass[conference,onecolumn]{IEEEtran}
\IEEEoverridecommandlockouts
\usepackage{cite}
\usepackage{amsmath,amssymb,amsfonts}
\usepackage{graphicx}
\usepackage{textcomp}
\usepackage{xcolor}
\usepackage{url}
\usepackage{hyperref}
\hypersetup{
    colorlinks=false,          % Must be false for the border to show
    pdfborderstyle={/S/U/W 1}, % S=Solid, U=Underline, W=Width (1pt)
    urlbordercolor=blue        % Color of the underline
}

\newcommand{\tbps}{Tb/s}
\newcommand{\pjbit}{pJ/bit}
 
\def\BibTeX{{\rm B\kern-.05em{\sc i\kern-.025em b}\kern-.08em
    T\kern-.1667em\lower.7ex\hbox{E}\kern-.125emX}}
 
\begin{document}
 
\title{Photonic Interconnects for TPU LLM Workloads}
 
\author{
\IEEEauthorblockN{Pranav Sitaraman}
\IEEEauthorblockA{sitaraman.pranav@gmail.com}
\and
\IEEEauthorblockN{Simone Surdi}
\IEEEauthorblockA{simonesurdi@gmail.com}
}
 
\maketitle
 
\noindent AccelForge with TCM (arXiv:2602.15172) finds optimal
per-chip mappings, but it has no concept of the network between
chips. We added one. The network model prices each transfer by hop
count, contention, multicast support, and a per-bit energy split
into a conversion term and a per-mm transmission term. An iterative
loop pushes that cost into a logical cluster-memory level so that
when we change the fabric the mapper retunes against
network-augmented memory. Above the mapper sits a discrete
distributed-schedule layer with four XLA/GSPMD-style policies, and
on top of that a stage model for a 9216-chip Ironwood pod. We
compare Lightmatter Passage (EVK50, EVK100, L200, L200X, M1000)
against an 800G electrical baseline and a strong 3.5~\tbps{},
6~\pjbit{} chiplet baseline on a 16-node cluster, a 64-chip
$4{\times}4{\times}4$ TPU cube, and the 9216-chip superpod, with six
LLM-stage GEMM proxies plus Gemma~4 and Gemini-class workloads.
 
 
\section{Key Insights and Innovations}
\begin{enumerate}
\item \textbf{Photonics doesn't move the inner mapping. It moves
the schedule.} Across 24 (workload, technology, scenario)
combinations on the 64-chip cube, AccelForge picked the same
partitioning, dataflow, and tile shape every time --- 0/24 inner
mappings changed. The wins live one layer up. The optimal
inter-chip communication pattern flipped on 5 of 16 workload-scenario
pairs, and on KV attention switching EVK100 from sharded collectives
to remote pooled-HBM streaming cut compute energy by 73.86\%.
Treating the per-chip mapping and the inter-chip schedule as one
number hides this, so we report them separately.
 
\item \textbf{The advantage isn't bandwidth but instead multicast and
distance-insensitive transport.} We picked the chiplet baseline
(3.5~\tbps{}, 6~\pjbit{}) specifically to rule out the
``photonics-is-just-more-bandwidth'' possibility. M1000 still beats it by
99.85\% on shared-DRAM multicast, because $N$ unicasts collapse to
one broadcast and contention drops to zero on a switched fabric.
Our link model separates conversion from transmission, so electrical
\pjbit{} grows with distance while photonic \pjbit{} stays flat. We find that per-chip multicast cost is generally flat from 64 to 9216 chips for M1000 while torus and chiplet fabrics face higher costs as $N$ grows.
 
\item \textbf{Photonics doesn't help everything equally.} On the
9216-chip Ironwood superpod, the M1000 stack delivers 42.43$\times$
inference speedup and 94.05\% energy reduction on the Gemma~4 26B
MoE proxy workload equivalent, while the long-context Gemini Deep Think proxy gets
14.50$\times$. Training is much smaller, 3.37$\times$ and
1.43$\times$ respectively. Once collectives shrink, training hits
an Amdahl ceiling at $\max(T_\text{compute}, T_\text{HBM})$ that
photonics doesn't move; inference stays exposed because decode and
KV/expert traffic cross the fabric every step. The strong case is
inference and long-context serving, not training.
\end{enumerate}
 
 
\section{Contributions}
\textbf{Simone Surdi.} Network model and AccelForge plumbing. Wrote
the topology classes (ring, mesh, 3-D torus, optical crossbar /
switched fabric) and the \texttt{LogicalClusterMemory} level.
Wrote \texttt{solve\_iteratively}, which re-runs the mapper with
updated logical-memory energies and latencies until they stop
moving. Built the mapping-signature comparison and ran the cluster,
TPU-cube, and chiplet-baseline studies. Pulled and verified the
technology constants against Lightmatter (SC25, ISES, product
pages) and TPU (Cloud docs, Ironwood blog, Hot Chips). Ran the
bandwidth and energy sensitivity sweeps. Helped write the report, poster, and technical summary.
 
\textbf{Pranav Sitaraman.} Distributed-schedule and superpod
analysis. Wrote the distance-parametrized link model and the OCS/OPS
sensitivity sweep. Implemented the four distributed schedules (SPMD
collective, remote pooled-HBM streaming, replicated cache, photonic
multicast cache fill). Built the 9216-chip Ironwood stage model
(prefill, decode, refill, MoE, training, gradient, optimizer) with
the $\alpha_\text{ovl}$ overlap factor. Ran strong and weak scaling
against the Gemma~4 anchors and Gemini-class proxies. Built the
ablation ladder and the 71-test integration suite. Helped write the report,
poster, and technical summary.
 
\smallskip
\noindent\emph{Joint.} Designed the four-layer pipeline (workload
$\to$ architecture $\to$ AccelForge mapping $\to$ distributed
schedule), worked through the experimental setup.
 
 
\section{Artifact}
\begin{itemize}
    \item \textbf{GitHub Repository:} \href{https://github.com/PranavSitaraman/photonic-interconnects}{https://github.com/PranavSitaraman/photonic-interconnects}
    \item \textbf{Report PDF:} \href{https://drive.google.com/file/d/1YP_ElsrpxuoX1Yr9KAZqQFVg7-6iGiRp/view?usp=sharing}{https://drive.google.com/file/d/1YP\_ElsrpxuoX1Yr9KAZqQFVg7-6iGiRp/view?usp=sharing}
    \item \textbf{Poster PDF:} \href{https://drive.google.com/file/d/1zd9CNyCml5L4fFaEEHRt--8OAQqAhu6B/view?usp=sharing}{https://drive.google.com/file/d/1zd9CNyCml5L4fFaEEHRt--8OAQqAhu6B/view?usp=sharing}
\end{itemize}
 
\end{document}